\documentclass[11pt,letterpaper]{article}

% Essential packages
\usepackage{fontspec}
% \usepackage[utf8]{inputenc}
% \usepackage[T1]{fontenc}
\usepackage{geometry}
\usepackage{amsmath}
\usepackage{amsfonts}
\usepackage{amssymb}
\usepackage{graphicx}
\usepackage{booktabs}
\usepackage{array}
\usepackage{longtable}
\usepackage{multirow}
\usepackage{multicol}
\usepackage{float}
\usepackage{caption}
\usepackage{subcaption}
\usepackage{hyperref}
\usepackage{fancyhdr}
\usepackage{datetime}
\usepackage{enumitem}
\usepackage{listings}
\usepackage{xcolor}
\usepackage{verbatim}
\usepackage{url}
\usepackage{svg}

% Page setup
\geometry{margin=1in}
\setlength{\parindent}{0pt}
\setlength{\parskip}{6pt}

% Header and footer setup
\pagestyle{fancy}
\fancyhf{}
\lhead{Data Science Capstone: Project Outline}
\rhead{\today}
\cfoot{\thepage}
\renewcommand{\headrulewidth}{0.4pt}

% Code listing setup
\lstset{
    basicstyle=\ttfamily\footnotesize,
    backgroundcolor=\color{gray!10},
    frame=single,
    breaklines=true,
    captionpos=b,
    numbers=left,
    numberstyle=\tiny\color{gray},
    keywordstyle=\color{blue},
    commentstyle=\color{green!60!black},
    stringstyle=\color{red}
}

% Custom commands
\newcommand{\heading}[1]{
    \section*{\workingtitle}
    \addcontentsline{toc}{section}{\workingtitle}
    \textbf{Date:} #1 \\
    \rule{\textwidth}{0.5pt}
}

\newcommand{\motivation}[1]{
    \subsection*{Motivation and Significance}
    \addcontentsline{toc}{subsection}{Motivation and Significance}
    #1
}

\newcommand{\plan}[1]{
    \section*{Project Plan and Timeline}
    \addcontentsline{toc}{section}{Timeline and Milestones}
    #1
}

\newcommand{\data}[2]{
    \section*{Datasets and Models}
    \addcontentsline{toc}{section}{Datasets and Models}
    \subsection*{Datasets}
    #1
    \subsection*{Models}
    #2
}

\newcommand{\problemDefBackground}[1]{
    \section*{Problem Definition and Background}
    \addcontentsline{toc}{section}{Problem Definition and Background}
    #1
}

\newcommand{\litReview}[1]{
    \section*{Literature Review}
    \addcontentsline{toc}{section}{Literature Review}
    #1
}

\newcommand{\methodology}[5]{
    \section*{Methodology}
    \addcontentsline{toc}{section}{Methodology}
    \subsection*{Overview}
    #1
    \subsection*{Models and Algorithms}
    #2
    \subsection*{Tech Stack}
    #3
    \subsection*{Implementation Strategy}
    #4
    \subsection*{Evaluation Metrics}
    #5
}




\newcommand{\mentor}[5]{
    \subsection*{Mentor Credentials}
    \addcontentsline{toc}{subsection}{Mentor Credentials}
    \begin{itemize}
        \item \textbf{Name:} #1
        \item \textbf{Email:} #2
        \item \textbf{Phone Number:} #3
        \item \textbf{Relevant Skills and Experience:} #4
        \item \textbf{LinkedIn:} #5
    \end{itemize}
}

% Title page setup
\newcommand{\workingtitle}{A Statistical Truth Detection System}
\title{\Large \textbf{DATS 598: Data Science Capstone} \\ 
       \large Project Outline \\
       \vspace{0.5cm}
       \normalsize \textbf{\workingtitle}}
\author{Eric Crisp \\ 
        University of Pennsylvania \\
        \texttt{ecrisp@upenn.edu}}
\date{\today}

\begin{document}

\maketitle
\tableofcontents
\newpage

% ========================================
% SAMPLE JOURNAL ENTRY -- MAIN PROPOSAL
% ========================================

% \heading{\today}


% =========================================
% summarize the project here !! 
% =========================================

\problemDefBackground{

This project is about taking input and determining if the statements within are true or false. Applications of this work include a variety of situations such as real-time news verification during debates, reports, or other media outlets, in researching the methods for retrieving evidence and verifying truth via RAG, embeddings, or other techniques, and finally the influence of truth in LLMs including response generation with and without truth verification.

The general result at the end of the day is a backend Python script that processes input and determines the truthfulness of the statements (and inferences) within. The backend will be wrapped around a simple web app for demonstration purposes, however, the main focus is the backend processing and truth verification. The application will be reflect a simple text prompt and will result in an analyzed output depending on the input and user configuration. 
}

% =========================================
% place the content for literatures here !! 
% =========================================

\litReview{ 

Summarize 4 to 5 research papers related to your problem, detailing approaches and outcomes.
How do these findings inform your project?


}

\methodology
{
% =========================================
% overall project methodlology goes here !! 
% =========================================
Overview of Approach: 
Explain the overall approach, including key steps and goals. Are you using an experimental, theoretical, or data-driven approach? Describe why you chose this approach.
Models and Algorithms: 
Briefly describe the specific models and algorithms you plan to use
Describe why these models are suitable for your problem and how they align with your project's objectives.
This subsection can brief, you will provide further detail in the next section
Tools and Software:
List the software, programming languages, or frameworks you intend to use (e.g., Python, R, TensorFlow, PyTorch, etc.).
Describe why these tools are appropriate and any benefits they offer to your project.
Implementation Strategy:
Outline how you plan to implement the models and algorithms. Will you be coding from scratch, or using existing libraries and resources?
Include information on data preprocessing, feature engineering, model training, validation, and testing.
Evaluation Metrics:
Describe how you will evaluate the success of your models. Discuss metrics such as accuracy, precision, recall, F1-score, etc.
Provide justification for your chosen metrics and their relevance to the problem at hand.

}
{
% =========================================
% what models and algos using here !! 
% =========================================
some stuff
}
{
% =========================================
% describe the tech stack details here !! 
% =========================================
some stuff
}
{
% =========================================
% implementation details here !! 
% =========================================
some stuff
}
{

% =========================================
% how to evaluate this here !! 
% =========================================
some stuff
}


\data{
% =========================================
% where and what data is being used !! 
% =========================================

All data being used is open source. FEVER, Factily, Wikidata, and any other data sources are (currently) being considered as "true" - verified and trusted sources.

\begin{itemize}

  \item \textbf{FEVER} – A large-scale dataset for fact-checking based on Wikipedia. \\
  \href{https://fever.ai}{https://fever.ai}

  \item \textbf{Factify} – A dataset and benchmark for factuality detection. \\
  \href{https://github.com/surya1701/Factify-2.0}{https://github.com/surya1701/Factify-2.0}

  \item \textbf{Wikidata} – A collaboratively edited knowledge base from the Wikimedia Foundation. \\
  \href{https://www.wikidata.org}{https://www.wikidata.org}

  \item \textbf{BERT} – A language model developed by Google AI. \\
  \href{https://arxiv.org/abs/1810.04805}{https://arxiv.org/abs/1810.04805}

  \item \textbf{T5} – A unified framework that casts all NLP tasks as text generation problems. \\
  \href{https://arxiv.org/abs/1910.10683}{https://arxiv.org/abs/1910.10683}

\end{itemize}
}
{
% =========================================
% where and what models are being used !! 
% =========================================

Identify and describe the datasets and models being considered.
Description of Datasets:
Provide an overview of the datasets you plan to use, including the source, size, structure, and format.
Discuss any preprocessing steps required to clean or transform the data, such as normalization, handling missing values, or encoding categorical variables.
Data Collection:
Describe how you will collect or source your datasets. Will you use public datasets, scrape data from websites, or collect it manually?
Highlight any ethical or legal considerations for data collection, especially if it involves sensitive or personal information.
Data Splitting and Training:
Explain how you will split the data into training, validation, and test sets. Justify the proportions and rationale behind this choice.
Describe any additional data augmentation or resampling techniques you plan to use.
Model Selection and Justification:
Detail the models or algorithms you will apply to the data, and explain why you chose them over other options. Mention any hyperparameters or specific settings you plan to explore for each model.
Potential Challenges and Mitigation:
Identify any potential challenges or risks associated with the datasets and models, such as class imbalance, overfitting, or underfitting.
Discuss how you plan to mitigate these risks and ensure robust, reliable results.



}

\plan{ 

Provide a rough timeline and list the updated project components.
Include soft deadlines for each significant phase.


\begin{enumerate}[label=\textbf{Week \arabic*:}]
    \setcounter{enumi}{3}

    \item \textbf{Create Initial Backend Model}
        \begin{itemize}[label=$\bullet$]
            \item Implement sentence embedding baseline using SentenceTransformers
            \item Build similarity-based classification
            \item Integrate Wikipedia API for evidence retrieval
            \item Create model evaluation framework
            \item \textbf{Deliverables:}
            \begin{itemize}[label=$\square$]
                \item Baseline fact-checker achieving on FEVER validation set
                \item Wikipedia API integration for real-time evidence retrieval
                \item Evaluation metrics (accuracy, F1, precision/recall)
                \item Create evidence ranking and filtering system
            \end{itemize}
        \end{itemize}

    \item \textbf{Build API Endpoints}
        \begin{itemize}
            \item Build hooks for external API calls in backend
            \item Create skeleton UX/UI features for displaying backend results
            \item \textbf{Deliverables:}
            \begin{itemize}
                \item Skeleton app and API endpoints in place.
            \end{itemize}
        \end{itemize}

    \item \textbf{Develop Backend Model Further}
        \begin{itemize}
            \item Implement advanced models (e.g., T5, BERT) via fine tuning on relevant datasets
            \item Experiment with RAG techniques for evidence retrieval (Pinecone)
        \end{itemize}

    \item \textbf{Determine and Test Performance Metrics}
        \begin{itemize}
            \item Sentence classification performance (F1).
            \item Document embedding and retrieval performance (MTEB)
            \item Truth detection performance (accuracy, precision/recall).
        \end{itemize}

    \item \textbf{Iterate on Backend Model}
        \begin{itemize}
            \item Implement sentence embedding baseline using SentenceTransformers
            \item Build similarity-based classification
            \item Integrate Wikipedia API for evidence retrieval
            \item Create model evaluation framework
        \end{itemize}

    \item \textbf{Create User Database}
        \begin{itemize}
            \item Implement a PostgreSQL database for user management (AWS RDS).
            \item Record user inputs and results for future analysis.
        \end{itemize}

    \item \textbf{Full Implementation and Testing of App}
        \begin{itemize}
            \item Finish API hooks and UX/UI features.
            \item Conduct end-to-end testing of the application.
            \item Create backlog of features to fix and improve.
        \end{itemize}

    \item \textbf{Conduct Initial "Search" Testing}
        \begin{itemize}
            \item Test output with and without truth detection to determine influence on LLM.
        \end{itemize}

    \item \textbf{Documentation and Final Report}
        \begin{itemize}
            \item Gather user documentation and create final report.
            \item Prepare presentation materials.
            \item Update Github README and documentation for public release.
        \end{itemize}


\end{enumerate}


}

\end{document}
